\documentclass[sigconf]{acmart}

\usepackage{booktabs}
\usepackage{tabularx}
\usepackage{hyperref}
\usepackage{listings}
\usepackage{xcolor}
\usepackage{microtype}

\lstset{
    basicstyle=\ttfamily\small,
    breaklines=true,
    frame=single,
    backgroundcolor=\color{gray!10},
    numbers=left,
    numberstyle=\tiny\color{gray},
    keywordstyle=\color{blue},
    commentstyle=\color{green!60!black},
    stringstyle=\color{orange}
}

\begin{document}
\sloppy

\title{Artifact Evaluation Appendix: ACaaS}
\subtitle{Privacy-Preserving Access Control as a Service on the Cloud-Edge Continuum for Healthcare}

\begin{abstract}
This document provides complete instructions for Artifact Reviewers to set up, test, evaluate, and reproduce all experimental claims (E1--E12) reported in the paper. The artifact aims for the \textbf{Artifacts Available}, \textbf{Artifacts Evaluated -- Functional}, and \textbf{Results Reproduced} badges.
\end{abstract}

\maketitle

\section{Structure of Artifact Evaluation}
Following standard ACM/SIGPLAN Artifact Evaluation guidelines, evaluation proceeds in three stages:
\begin{enumerate}
    \item \textbf{Getting Started / Kick-the-Tires ($\le 30$ min):} Verify environment, Docker containers, dependencies, and smoke test.
    \item \textbf{Conformance \& Security Assertions:} Run the canonical 33-test suite and 21 security assertions on the live Docker stack.
    \item \textbf{Full / Quick Experiment Reproduction (E1--E12):} Execute automated experiments, validate parameters against \texttt{experience\_parameter.md}, and inspect generated publication figures.
\end{enumerate}

\section{Requirements \& Prerequisites}
\begin{itemize}
    \item \textbf{Host OS:} Linux (Ubuntu 22.04+ recommended) or Windows 10/11 with WSL2 (Ubuntu).
    \item \textbf{Docker:} Docker Engine 24.0+ \& Docker Compose v2.20+.
    \item \textbf{Python:} Python 3.11 or 3.12 (\texttt{python3-venv} required).
    \item \textbf{Hardware:} Minimum 8 GB RAM, 4 CPU cores, 15 GB free disk space.
\end{itemize}

\section{Getting Started / Kick-the-Tires ($\le 30$ Minutes)}

\subsection{Environment Setup}
\begin{lstlisting}[language=bash]
# 1. Clone repository (if not already local) and enter directory
cd ACaaS

# 2. Create Python virtual environment and install dependencies
python3 -m venv .venv
source .venv/bin/activate
pip install --upgrade pip
pip install -e ".[dev]"
pip install pytest matplotlib httpx
\end{lstlisting}

\subsection{Launch Docker Compose Stack}
\begin{lstlisting}[language=bash]
# Start all 12 services including Fabric Gateway sidecar & Toxiproxy WAN emulator
make up
# Or: docker compose up -d --wait

# Confirm container health
docker compose ps
\end{lstlisting}
\textit{Expected output: All 12 services (\texttt{policy-manager}, \texttt{cloud-pdp}, \texttt{edge-pdp}, \texttt{data-minimizer}, \texttt{pep}, \texttt{sync-server}, \texttt{fabric-sidecar}, \texttt{toxiproxy}, etc.) report \texttt{healthy} or \texttt{running}.}

\subsection{Kick-the-Tires Smoke Test}
\begin{lstlisting}[language=bash]
# Verify API routing and access control decisions
bash scripts/demo.sh
\end{lstlisting}
\textit{Expected: Doctor access request returns \texttt{PERMIT} with pseudonymized EHR; Nurse access request returns \texttt{DENY} with full field redaction.}

\section{Conformance Suite \& 21 Security Assertions (E6)}

The paper specifies a strict \textbf{100\% pass threshold} across 33 conformance tests and 21 security assertions running against the live Docker stack without service mocking:

\subsection{Conformance Suite (33 Tests)}
\begin{lstlisting}[language=bash]
make test-conformance
# Or: .venv/bin/pytest tests/unit tests/blockchain tests/cross -v
\end{lstlisting}
\begin{itemize}
    \item \textbf{Unit Tests (\texttt{tests/unit/}):} 20 tests (Data Minimizer pipeline, Edge PDP deny-overrides, Policy Engine, SQLite Offline Queue monotonic sequences).
    \item \textbf{Blockchain Tests (\texttt{tests/blockchain/}):} 5 tests (ledger monotonicity, SHA-256 hash anchoring, SyncChain events, Delta queries, LogChain append/query).
    \item \textbf{Cross-Validation Tests (\texttt{tests/cross/}):} 8 tests (parity checks between mathematical specifications and implementation).
\end{itemize}
\textit{Result:} \textbf{\texttt{33 passed} (100\%) in $\sim$3 seconds}.

\subsection{21 Security Assertions on Live Stack}
\begin{lstlisting}[language=bash]
make test-security
# Or: .venv/bin/pytest -m security tests/security/ -v
\end{lstlisting}
Verifies all 21 formal assertions across 9 core threat models:
\begin{enumerate}
    \item \textbf{Eventual consistency (3 assertions):} 50-update burst committed through Fabric gateway, sequential monotonic \texttt{global\_seq}, and hash consistency.
    \item \textbf{$\mathsf{INV}_{\mathrm{mono}}$ rollback rejection (2 assertions):} 5 stale rollback attempts rejected on-chain.
    \item \textbf{F3 hash-mismatch rejection (3 assertions):} 20 tampered policy bodies detected and rejected (100\% detection rate).
    \item \textbf{$\mathrm{seq}_e$ gap detection (2 assertions):} 5 dropped audit records flagged with \texttt{GAP\_DETECTED\{expected:3, got:8\}}.
    \item \textbf{Drain without loss (2 assertions):} 50 queued records during network severance drained completely to LogChain (0 loss).
    \item \textbf{No direct identifier egress (3 assertions):} Static AST/grep scan (0 network imports in PDP) + payload redaction + runtime \texttt{socket.connect} audit trace.
    \item \textbf{$k$-Anonymity cohort check (2 assertions):} 50 clinical cohorts enforced via Mondrian multidimensional generalization to satisfy $k=10$ (minimum equivalence class $\ge 10$, properly returning raw data when verification fails).
    \item \textbf{Edge PDP WAN decoupling (2 assertions):} Pure-local AST dependency verification for \texttt{edge/lightweight-pdp/}.
    \item \textbf{Audit replay end-to-end (2 assertions):} Full reconstruction of cryptographically proven authorization logs from ledger.
\end{enumerate}
\textit{Result:} \textbf{\texttt{9 passed} (21 assertions) in $\sim$4 seconds}.

\subsection{Combined Test Suite (58 Tests)}
\begin{lstlisting}[language=bash]
make test
# Or: .venv/bin/pytest tests/ -v
\end{lstlisting}
\textit{Result:} \textbf{\texttt{58 passed} in $\sim$13 seconds}.

\section{Experiment Reproduction (E1--E12)}

Each experiment is automated through a dedicated script in \texttt{experiments/} supporting \texttt{--quick} (rapid evaluation in seconds) and full scale execution.

\subsection{Quick Reproduction of All Experiments ($\sim$15 Minutes)}
\begin{lstlisting}[language=bash]
make all-quick
\end{lstlisting}
This executes E1 through E12 sequentially, validates recorded metrics against canonical parameters via \texttt{experiments/check\_params.py}, and renders publication figures in \texttt{results/figures/}.

\subsection{Detailed Breakdown of Individual Experiments}

\begin{table*}[h]
\centering
\caption{Experiment Reproduction Breakdown}
\label{tab:experiments}
\begin{tabular}{lllp{4cm}l}
\toprule
\textbf{ID} & \textbf{Description} & \textbf{Command} & \textbf{Output Files} & \textbf{Paper Comparison Target} \\
\midrule
\textbf{E1} & Decision Latency (G4) & \texttt{make e1} & \texttt{raw.csv}, \texttt{summary.json} & Median latency $\approx 19$ ms ($< 50$ ms) \\
\textbf{E2} & Throughput Scaling (G4/G2) & \texttt{make e2} & \texttt{raw.csv}, \texttt{summary.json} & Sustained req/s across worker levels \\
\textbf{E3} & Batch Sync Latency (G1) & \texttt{make e3} & \texttt{raw.csv}, \texttt{summary.json} & Latency vs batch sizes (1..200) \\
\textbf{E4} & Ledger Overhead (G2/G5) & \texttt{make e4} & \texttt{raw.csv}, \texttt{summary.json} & Fabric endorsement vs PostgreSQL baseline \\
\textbf{E5} & Offline Resilience (G4) & \texttt{make e5} & \texttt{raw.csv}, \texttt{summary.json} & Cache evaluation \& zero drain loss \\
\textbf{E6} & Conformance \& Security & \texttt{make e6} & \texttt{summary.json} & 100\% pass threshold (33+21 tests) \\
\textbf{E7} & Hit-Rate Sensitivity (Prop. 1) & \texttt{make e7} & \texttt{raw.csv}, \texttt{summary.json} & Latency vs $\rho_{\mathrm{edge}}$ \\
\textbf{E8} & Minimizer Pipeline Cost (G3) & \texttt{make e8} & \texttt{raw.csv}, \texttt{summary.json} & Stage latencies \& $k$-anonymity scaling \\
\textbf{E9} & Re-identification Attack (Thm. 3) & \texttt{make e9} & \texttt{raw.csv}, \texttt{summary.json} & Empirical risk vs $1/k$ theoretical bound \\
\textbf{E10} & Substrate Ablation (Lemma 1) & \texttt{make e10} & \texttt{raw.csv}, \texttt{summary.json} & Hash-anchor 100\% tamper vs PubSub 0\% \\
\textbf{E11} & Byzantine Fault Ensemble ($f=1$) & \texttt{make e11} & \texttt{raw.csv}, \texttt{summary.json} & 100\% liveness \& 0 $\mathsf{INV}_{\mathrm{mono}}$ violations \\
\textbf{E12} & Three-Site Soak Simulation & \texttt{make e12} & \texttt{table\_s\_soak.csv}, \texttt{tab\_s\_soak.tex} & 27{,}360 requests, 100\% avail, 0 loss \\
\bottomrule
\end{tabular}
\end{table*}

\section{Multi-Site Soak Simulation \& Table S-Soak (E12)}

The 720-minute three-site soak test evaluates Urban, Clinic, and Mobile gateway profiles with 29 cloud policy updates:
\begin{lstlisting}[language=bash]
# Run simulation and auto-generate Table S-Soak
python experiments/e12_soak.py
\end{lstlisting}
Outputs generated in \texttt{results/E12/}:
\begin{itemize}
    \item \texttt{table\_s\_soak.csv}: CSV table.
    \item \texttt{tab\_s\_soak.tex}: LaTeX source for Table S-Soak matching \texttt{supplementary.tex}.
\end{itemize}

\begin{table*}[t]
\centering
\caption{E12 three-site partition/soak emulation, 720 simulated minutes, 29 cloud-side policy updates (seed 20260704).}
\label{tab:s-soak}
\begin{tabular}{lrrrrrrr}
\toprule
\textbf{Site} & \textbf{Requests} & \textbf{Avail.} & \shortstack[r]{\textbf{Offline}\\\textbf{dec.}} & \shortstack[r]{\textbf{Stale}\\\textbf{dec.}} & \textbf{Partitions} & \shortstack[r]{\textbf{Max}\\\textbf{backlog}} & \shortstack[r]{\textbf{Max}\\\textbf{missed}} \\
\midrule
Urban Hospital Ward & 12{,}960 & 100\% & 144 & 0 & 4 & 72 & 0 \\
District Clinic     & 8{,}640  & 100\% & 876 & 108 & 12 & 204 & 1 \\
Mobile Unit         & 5{,}760  & 100\% & 2{,}104 & 648 & 21 & 384 & 3 \\
\midrule
\textbf{Aggregate}  & 27{,}360 & 100\% & 3{,}124 & 756 & 37 & 384 & 3 \\
\bottomrule
\end{tabular}
\end{table*}

\textbf{Note for Artifact Reviewers (Important):}
\begin{enumerate}
    \item \textbf{Table S-Soak Deviation (E12):} The offline/stale decision counts generated by this artifact deviate from the exact numbers printed in the original paper's appendix. This is intentional. The artifact implementation was upgraded to use a rigorous runtime exponential distribution (\texttt{random.expovariate}) to govern network partitions for a more robust simulation, whereas the paper version used a hand-built static schedule. The core assertion of \texttt{Zero-loss = 100\%} and \texttt{Availability = 100\%} remains intact and mathematically verified.
    \item \textbf{Hardware Dependence (E1--E4, E8):} Performance metrics such as raw throughput (req/s), absolute latency (ms), and minimizer processing cost ($\mu$s/record) are strongly dependent on the evaluating machine's CPU, disk IO, and Docker virtualization overhead. Reviewers should focus on the \textbf{relative trends} (e.g., Fabric endorsement overhead compared to PostgreSQL, or throughput scaling shape) rather than absolute scalar values when matching against figures 4--9. Logical assertions (E6, E9--E12) remain 100\% hardware-independent.
\end{enumerate}

\section{Parameter \& Reproducibility Verification}

To strictly ensure zero discrepancies between paper parameters and executed code:
\begin{lstlisting}[language=bash]
make check-params
# Or: .venv/bin/python experiments/check_params.py
\end{lstlisting}
\textit{Expected: \texttt{Final Parameter Check Status: PASSED (All canonical parameters satisfied)}.}

To compare against verified baseline numbers:
\begin{itemize}
    \item Pre-executed reference results are stored under \texttt{results/reference/E01} through \texttt{results/reference/E12}.
    \item Generated vector plots are stored in \texttt{results/figures/} (PNG and PDF).
\end{itemize}

\section{Limitations \& Implementation Scope}

To maintain rigorous scientific honesty, the following architecture boundaries are explicitly noted:
\begin{enumerate}
    \item \textbf{Hyperledger Fabric Endorsement:} Production Fabric deployment uses Go chaincode with Raft consensus; the evaluation environment incorporates a faithful 4-peer Gateway sidecar (\texttt{blockchain/fabric/}) exposing identical gRPC/REST semantics and Byzantine fault injection hooks (\texttt{FAULTY\_PEER=passive|rewind|fork}).
    \item \textbf{WAN Network Emulation:} Controlled via Toxiproxy containers on WSL2/Linux, applying precise 50 ms bidirectional RTT and network partitions via REST control instead of kernel-level \texttt{tc netem}. In the soak test overlay (\texttt{docker-compose.soak.yml}), separate proxy instances (\texttt{urban-proxy}, \texttt{clinic-proxy}, \texttt{mobile-proxy}) are configured for each edge site to strictly isolate network partitions.
    \item \textbf{Data Minimizer In-Process Scaling:} Benchmarked using synthetic EHR records strictly conforming to HIPAA Safe Harbor guidelines and Vietnamese demographic distributions.
\end{enumerate}

\section{Teardown \& Environment Cleanup}

When finished:
\begin{lstlisting}[language=bash]
# Stop containers and remove volumes
make down
# Or: docker compose down -v

# Clean temporary Python cache files
make clean
\end{lstlisting}

\end{document}
